\section{Extended Observation: Bearing and Standoff From One Crossing}\label{sec:observation}

A validated crossing gives a point on the cable, and the same records also give the cable bearing and the standoff. The point $\mathbf{p}_c$ is the midpoint of the two sensors' fixes, or the vehicle position at the GZ fix. The positions of the two sensors are computed from the navigation solution:
\begin{equation}
\mathbf{p}_{UNO,k} = \mathbf{p}_k - \frac{b}{2}\,\hat{\boldsymbol{\ell}}(\chi_k), \qquad \mathbf{p}_{DUO,k} = \mathbf{p}_k + \frac{b}{2}\,\hat{\boldsymbol{\ell}}(\chi_k), \qquad \hat{\boldsymbol{\ell}}(\chi) = (\cos\chi, -\sin\chi).
\label{eq:sensor_positions}
\end{equation}
Here $\mathbf{p}_k$ is the position (East, North) given by the global navigation satellite system (GNSS) at sample $k$, $\chi_k$ the yaw given by the attitude and heading reference system (AHRS) and $\hat{\boldsymbol{\ell}}$ the starboard unit vector.

The classical closed-form estimators of bearing and standoff read the field at individual samples, as a ratio of components or a peak value. At low signal-to-noise ratio (SNR) a single sample is noisy. TDOD, RAT, GRAT and AMP below therefore work on projections of the records onto the cable OBF $\{\Phi_1, \Phi_2\}$ or the gradiometer OBF $\{\Gamma_1, \Gamma_2\}$ of \cref{sec:setup}, which integrate the whole crossing. TRI is kept in its sample-wise form as the reference method of the cable-tracking literature. All of them are in closed form. Constants are treated as nuisance parameters throughout, through a constant column in each least-squares fit, as in the matched-subspace detector with interference~\cite{scharf1994matched}. \Cref{tab:observers} lists the estimators, what each one needs, and the raw-field method it replaces.

\begin{table}[htb]
\centering
\small
\caption{Bearing and standoff estimators of one crossing. $S$ is the total field, Attitude is what the estimator needs from the AHRS, and Form is a closed form, a same-shape ratio or a closed-form step iterated to a fixed point. AMP and TRI also use the bearing, through the normal coordinate~\eqref{eq:normal_coordinate} or the incidence $\theta$.}
\label{tab:observers}
\begin{tabular}{@{}llllll@{}}
\toprule
Estimator & Channel & Attitude & Current & Form & Raw-field method \\
\midrule
TDOD & ZC fixes of $B_z$ & full & no & closed & line through two fixes~\cite{bharti2020semi} \\
RAT & $B_E$, $B_N$ & full & no & ratio & heading deviation~\cite{xiang2016subsea} \\
GRAT & $G_E$, $G_N$ & heading & no & ratio & gradient direction~\cite{kdd1998maritime} \\
FUS & TDOD, RAT, GRAT & as inputs & no & closed & -- \\
\midrule
AMP & $B_z$ or $S$ & full or none & yes & fixed point & level, known current~\cite{kdd1998maritime} \\
TRI & $B_n$, $B_z$, two sensors & full & no & closed & two inclinations~\cite{xiang2016subsea} \\
\bottomrule
\end{tabular}
\end{table}

\subsection{Bearing}\label{subsec:bearing}

\paragraph{Time difference of detection (TDOD).}
Both sensor fixes lie on the cable, so the bearing is the direction that joins them, as in crossing-based trackers~\cite{bharti2020semi}:
\begin{equation}
\hat{\varphi}_{TDOD} = \operatorname{atan2}\left(E_{DUO} - E_{UNO},\, N_{DUO} - N_{UNO}\right) \bmod 180^\circ.
\label{eq:tdod}
\end{equation}
$(E, N)$ are the East and North coordinates of each fix in metres, and $\varphi$ is measured clockwise from north. Because it is written in positions, TDOD absorbs speed variations. Along-track errors $\varepsilon_1$ and $\varepsilon_2$ of the UNO and DUO fixes give the bearing error
\begin{equation}
\Delta\hat{\varphi}_{TDOD} \simeq \frac{(\varepsilon_2 - \varepsilon_1)\sin^2\theta}{b},
\label{eq:tdod_error}
\end{equation}
which grows with the standoff, since the fix errors scale with $d/\sin\theta$. The two fixes are the ZC fixes of $B_z$, so TDOD depends on the attitude. GZ gives a single fix, so TDOD does not apply to it.

\paragraph{Same-shape ratios.}
The two other bearings rest on one property. When two channels have the same waveform up to a constant factor, the ratio of their least-squares coefficients on any template equals that factor exactly, whatever the template scale, its position and the truncation of the record (Appendix~\ref{app:same_shape}). The template only sets the SNR of the result, which is highest when the template matches the waveform.

\paragraph{OBF ratio (RAT).}
The vector channels are rotated to the North--East--Down (NED) frame with the AHRS attitude. In NED the horizontal anomaly vector is aligned with the cable normal at every sample and follows the even profile $\Phi_1$ (\cref{tab:channels}). Its East and North components, $B_n \sin(\varphi + 90^\circ)$ and $B_n \cos(\varphi + 90^\circ)$, therefore have the same shape, and their ratio gives the normal. This is the heading-deviation relation of~\cite{xiang2016subsea,zhang2024electromagnetic}, written on projections rather than samples:
\begin{equation}
\hat{\varphi}_{RAT} = \operatorname{atan2}(Q_E, Q_N) + 90^\circ \bmod 180^\circ,
\label{eq:rat}
\end{equation}
where $Q_j$ is the matched-filter output of $B_j$ at the fix lag, summed over the two sensors,
\begin{equation}
Q_j = \sum_{s \in \{UNO, DUO\}} \langle B_{j,s}, \tilde{\Phi}_{1,s} \rangle, \qquad j \in \{E, N\},
\label{eq:rat_projection}
\end{equation}
and $\tilde{\Phi}_{1,s}$ is the unit-norm template $\Phi_1$ centred on the fix of sensor $s$ and made zero-mean over its window, which rejects the offsets of the components (Appendix~\ref{app:same_shape}). As a same-shape ratio, RAT is exact under truncation and for any template scale, but it needs the full attitude. On the asymmetric windows of a snake leg the slow attitude drift does not cancel and leaks into the bearing.

\paragraph{Gradiometer ratio (GRAT).}
The horizontal gradiometer vector $(G_E, G_N)$, the horizontal part of $\mathbf{B}_{DUO} - \mathbf{B}_{UNO}$, also lies along the cable normal, with the odd profile $\Gamma_1$~\eqref{eq:grad_obf}. Its two components are again a same-shape pair, and $Q^G_j$ is the matched-filter output of $G_j$ with the template $\Gamma_1$ centred on the fix and made zero-mean over its window:
\begin{equation}
\hat{\varphi}_{GRAT} = \operatorname{atan2}(Q^G_E, Q^G_N) + 90^\circ \bmod 180^\circ.
\label{eq:grat}
\end{equation}
This is the direction of the field gradient used with differential sensors~\cite{kdd1998maritime,zeng2023magnetic}, written on projections. Common-mode fields cancel in the difference, tilt errors enter at second order, and only the heading is needed. The amplitude is proportional to $|\cos\theta|$ and vanishes at perpendicular incidence.

\paragraph{Fusion (FUS).}
Bearings are axial, defined modulo $180^\circ$, so they are fused by an inverse-variance mean of the doubled angle~\cite{mardia2000directional}:
\begin{equation}
\hat{\varphi}_{FUS} = \frac{1}{2} \operatorname{atan2}\left(\sum_i w_i \sin 2\hat{\varphi}_i,\, \sum_i w_i \cos 2\hat{\varphi}_i\right), \qquad w_i = \sigma_i^{-2},
\label{eq:fus}
\end{equation}
with $i \in \{\mathrm{TDOD}, \mathrm{RAT}, \mathrm{GRAT}\}$. The $\sigma_i$ are trimmed RMS errors read from tables indexed by standoff and incidence, calibrated on simulated single crossings \todo{(Section 6, E1)}. The weights assume independent errors, which holds only approximately since TDOD and RAT share the AHRS. FUS combines the three bearings when the attitude is trusted, and GRAT is used alone otherwise.


\subsection{Standoff}\label{subsec:standoff}

The coordinate of each sample along the estimated cable normal is
\begin{equation}
x_k = (\mathbf{p}_{s,k} - \mathbf{p}_c) \cdot \hat{\mathbf{n}}(\hat{\varphi}), \qquad \hat{\mathbf{n}}(\varphi) = (\cos\varphi, -\sin\varphi),
\label{eq:normal_coordinate}
\end{equation}
where $\mathbf{p}_{s,k}$ is the position~\eqref{eq:sensor_positions} of sensor $s$. A bearing error $\Delta\varphi$ rescales $x$ by about $(1 - \Delta\varphi \cot\theta)$ (\cref{app:amp}).

\paragraph{Amplitude with known current (AMP).}
The feed current of a repeatered cable is regulated by the shore station and known to the operator~\cite{chesnoy2016undersea}, and cable-locating systems have long used it to convert a field level into a distance~\cite{kdd1998maritime}. In the OBF the conversion uses a single coefficient, that of the channel on the odd function $\Phi_2$. For the vertical component this coefficient is $A_z = \sqrt{\pi/2}\,K_p$, and for the total field it is $A_2 = \sqrt{\pi/2}\,K_p \sin\delta$~\eqref{eq:total_field}, which does not depend on the cable bearing. With $K_p = \mu_0 I/(2\pi d)$ the standoff follows from the coefficient $\hat{A}$ of either channel,
\begin{equation}
\hat{d} = \frac{\mu_0 I}{2\pi}\,\sqrt{\frac{\pi}{2}}\,\frac{g}{|\hat{A}|}, \qquad g = 1 \text{ for } B_z, \qquad g = \sin\delta \text{ for the total field}.
\label{eq:amp}
\end{equation}
$\hat{A}$ is the $\Phi_2$ coefficient of the least-squares fit of $[\Phi_1(x_k/d), \Phi_2(x_k/d), 1]$ to the channel on a window of $\pm K$ CPA distances around the fix, averaged over the two sensors. Because this coefficient depends on the template scale $d$, AMP applies~\eqref{eq:amp} as a fixed-point iteration,
\begin{equation}
d_{j+1} = \frac{\mu_0 I}{2\pi}\,\sqrt{\frac{\pi}{2}}\,\frac{g}{|\hat{A}(d_j)|},
\label{eq:amp_iteration}
\end{equation}
started from the prior standoff $d_0$ and stopped when the change falls below 1~cm. The true standoff is a fixed point for every truncation of the record, because a template at the right scale fits the record exactly. Each step halves the error on a long pass and reduces it much faster on a snake leg, so a few iterations suffice and the result does not depend on the starting value (\cref{app:amp}). Each step is a linear least-squares fit, so AMP needs no scan. With a trusted attitude AMP can use $B_z$. Otherwise it uses the total field, which needs no attitude, and together with GZ and GRAT the whole observation then needs only the heading. The bearing enters only through the template scale, and on a snake leg its effect is a small fraction of its effect on a fit of the profile shape (\cref{app:amp}).

The price of AMP is a direct sensitivity to the assumed current. The effective magnetic current can differ from the shore-station current, for instance through return currents along the armour near shore ends, and a current error of 10\% biases $\hat{d}$ by about 10\%. Heave modulates the amplitude in the same way.

\paragraph{Two-sensor triangulation (TRI).}
When the cable lies between the sensors ($x_{UNO} x_{DUO} < 0$), each sensor's inclination gives $|B_{z,i}|/|B_{n,i}| = |x_i|/d$, and $|x_{UNO}| + |x_{DUO}| = b|\cos\theta|$~\cite{xiang2016subsea}, hence
\begin{equation}
\hat{d}_{TRI} = \frac{b|\cos\theta|}{|B_{z,UNO}|/|B_{n,UNO}| + |B_{z,DUO}|/|B_{n,DUO}|},
\label{eq:tri}
\end{equation}
taken as the median over the samples where the cable lies between the sensors. TRI is the reference two-sensor method of the cable-tracking literature~\cite{xiang2016subsea,yu2018robust,zhang2024electromagnetic}, and it needs no current. It measures $d$ against the lateral separation of the sensors, $b|\cos\theta|$, which is several times smaller than the standoff here, so small errors on the inclinations become large errors on $d$. The trivial baseline, the prior mean of 3~m, is reported alongside.


\subsection{Observability}\label{subsec:observability}

With position fixes only, one crossing constrains the cable to an infinite set of lines through a point. Two crossings are then needed for the bearing, and the standoff is never observed~\cite{bharti2020semi}. The bearing observation makes the horizontal line observable from a single crossing, and the standoff observation adds the depth of the cable below the sensors. Their practical value depends on their accuracy compared with the variability of the cable along the route, which \todo{Section 7} quantifies.

